% SOURCE: 项目说明书-初稿.md §七 + 整体总结
\section{结论}

本作品以双摆这一经典非线性动力学系统为研究对象，构建了一个集高精度数值仿真、丰富可视化手段、
完整教学闭环于一体的虚拟仿真实验平台。本报告从物理原理、系统架构、数值方法、使用方法与结果验证五个维度进行了系统论述，
主要成果总结如下：

\begin{enumerate}[label=(\arabic*), leftmargin=*]
  \item \textbf{物理原理严格正确}：运动方程由拉格朗日力学严格导出，通过小角度线性化（偏差 $< 0.8\%$）、
       单摆退化（周期偏差 $< 0.05\%$）、能量守恒（漂移 $< 0.38\%$）三项物理验证，
       为仿真系统的可信度提供了可复现的量化背书；

  \item \textbf{数值方法自研可控}：从 ODE 右端函数到 RKF45 积分器、自适应步长控制、Web Worker 并行管线、
       Transferable 零拷贝传输，全部核心算法均自主实现，无第三方物理引擎依赖，
       为教学中的算法透明性提供了保障；

  \item \textbf{教学架构设计合理}：四层递进认知闭环（感知 $\to$ 分析 $\to$ 验证 $\to$ 创造）将
       混沌教学从“现象演示”升级为完整的探究式学习路径，覆盖从本科课堂到科研训练的完整需求梯度。
       蝴蝶效应双视口对比器与时间反演双模式实验作为两个标志性功能，
       将“初值敏感性”和“决定论 vs 可预测性”这两个混沌教学的核心难点转化为直观可交互的视觉叙事；

  \item \textbf{部署门槛极低}：纯客户端架构（零后端依赖、零网络需求），
       产出物为自包含 HTML5 文件夹——双击即可运行，可嵌入课程网站、拷入 U 盘、部署在教室一体机或科技馆终端，
       具备大规模推广的技术条件。
\end{enumerate}

\subsection*{运行配置要求}

\begin{table}[htbp]
  \centering
  \caption{系统运行环境要求}
  \label{tab:requirements}
  \begin{tabularx}{\textwidth}{l l}
    \toprule
    \textbf{项目} & \textbf{要求} \\
    \midrule
    浏览器 & Chrome 90+ / Firefox 90+ / Edge 90+ / Safari 15+ \\
    操作系统 & Windows 10 及以上 / macOS 11 及以上 / Linux（主流发行版） \\
    内存 & 推荐 4 GB 及以上 \\
    显卡 & 支持 WebGL 2.0 的显卡（集成显卡可运行，独立显卡推荐） \\
    网络 & 离线可用；Python 沙箱功能首次使用需在线加载 Pyodide（约 18 MB） \\
    \bottomrule
  \end{tabularx}
\end{table}

\subsection*{访问方式}

本作品以自包含 HTML5 离线文件夹形式提供，双击 \texttt{index.html} 即可在浏览器中运行，
无需网络连接或本地服务器，详见竞赛报名材料。

\endinput
